\chapter[Phase 10 corrections]{Phase 10: corrections to the results of Phases 0--5}\label{ch:cIII6}
\chapterbrief
{the setting of Phases~1--3: the commuting-terms class $\mathfrak B$ on a bounded rectangle (Chapter~\ref{ch:cI3}), and the factor decomposition $\ker P=\sum\ker Q_j$ in Sobolev spaces on bounded Lipschitz regions (Chapter~\ref{ch:cI2}).}
{to record the proofs that changed what Phases~0--5 had claimed: a statement that turned out false as written (the continuity no-go for any operator with two terms), its correct form, and the counterexample showing that density of the sum of kernels fails on regions with holes, which also settles the question ``closed but not dense''.}
{Dim.~A52--A53. They prove Chapter~\ref{PAP-ch:phase10} of the paper. They prevail over the statements of Phases~0--5 with the same identifiers (ledger sections F10-A, F10-C2, F10-TRIAGE).}

\section{The continuity no-go with constant terms}

\begin{notebox}{Setting}
$\Omega=I\times J$ is a bounded rectangle, $I=(0,T)$, $\mathfrak B$ the class of Chapter~\ref{ch:cI3}, $t_0\in\bar I$ the base time, $s=t-t_0$. Consider
\[
Q=\tau-\sum_{n=1}^Na_n(\xi)-\sum_{k=1}^Kc_k,\qquad a_n\in\C[\xi]\ \text{non-constant},\ c_k\in\C,
\]
$a=\sum a_n$, $\kappa=\sum c_k$, $Q'=\tau-a(\xi)$. A \emph{term} is a monomial $\lambda_k\xi^{l_k}$; a monomial of degree $0$ (linear reaction $\lambda_0\,\mathrm{id}$) is a term. On $\mathfrak B\cap\ker Q$,
\[
E^{\rm add}(u)(t)=u(t)-u_0-\sum_n\big(e^{sa_n(\partial)}u_0-u_0\big)-\sum_k\big(e^{sc_k}-1\big)u_0,
\]
and $E'$ is the same object for $Q'$.
\end{notebox}

\begin{dimcard}{A52}{S10A-T1}{P1-Q6: the no-go for $\Eint$ holds with two non-constant terms, and fails with constants only (S10A-L1, T1, C1)}{\proved}{thm:s10at1}
\begin{enumerate}[label=(\alph*)]
\item (S10A-L1) $M:v\mapsto e^{\kappa(t-t_0)}v$ is a linear bijection $\mathfrak B\cap\ker Q'\to\mathfrak B\cap\ker Q$ with $c_m\|v\|_{H^m}\le\|Mv\|_{H^m}\le C_m\|v\|_{H^m}$, and for $u=Mv$
\[
E^{\rm add}(u)=E'(v)+B(v),\qquad B(v)=\big(e^{s\kappa}-1\big)v(t)-\sum_k\big(e^{sc_k}-1\big)v_0,\quad \|B(v)\|_{L^2}\le C\|v\|_{H^1}.
\]
\item (S10A-T1) For $m\ge1$, $E^{\rm add}$ is bounded $H^m\to L^2$ iff $E'$ is; if $E'$ is unbounded in $H^m$ for every $m\ge0$, so is $E^{\rm add}$; if $N\le1$, $E^{\rm add}$ is bounded for every $m\ge1$.
\item (S10A-C1) For monomial terms of distinct degrees $0\le l_1<\dots<l_M$: $E^{\rm add}$ is unbounded $H^m\to L^2$ for all $m$ and $t_0$ iff $\#\{k:l_k\ge1\}\ge2$; otherwise bounded for every $m\ge1$.
\end{enumerate}
Consequently the literal statement of P1-Q6, ``no continuity for any operator with at least two non-proportional terms'', is \textbf{false}: reaction--diffusion $u_t=Du_{xx}+ru$ has $E^{\rm add}$ continuous from $H^1$, with
\[
E^{\rm add}(u)=(e^{sr}-1)\big(v(t)-v_0\big),\qquad G_\zeta(s)=(e^{sD\zeta^2}-1)(e^{sr}-1)\ \text{on modes }e^{a(\zeta)t+\zeta x}.
\]
\end{dimcard}

\begin{proofidea}
A constant term only multiplies the solution by $e^{\kappa s}$ and shifts the interaction by a bounded, lower-order piece; so it neither creates nor destroys the no-go, which comes from the non-constant terms.
\end{proofidea}

\begin{proof}
\step{1}{(a) The constants separate}
If $u=e^{\kappa s}v$ then $\partial_tu-a(\partial_x)u-\kappa u=e^{\kappa s}(\partial_tv-a(\partial_x)v)$, so $Q(\partial)u=0\iff Q'(\partial)v=0$. The class $\mathfrak B$ is stable under multiplication by $e^{\gamma t}$: by Leibniz $e^{\gamma t}\mathfrak B_\rho\subset\mathfrak B_{\rho+|\gamma|}$, and the $H^m$ estimates follow with $e^{\pm\kappa s}$ smooth and bounded with all derivatives on $\bar I$. With $u_0=v_0$ and $u(t)=e^{\kappa s}v(t)$, substituting in the definition gives $E^{\rm add}(u)-E'(v)=(e^{\kappa s}-1)v(t)-\sum_k(e^{sc_k}-1)v_0=B(v)$. The bound $\|B(v)\|_{L^2(\Omega)}\le C\|v\|_{H^1(\Omega)}$ uses the trace $\|v_0\|_{L^2(J)}\le C_I\|v\|_{H^1(\Omega)}$.

\step{2}{(b) Transfer of boundedness}
$E^{\rm add}\circ M=E'+B$ with $B$ bounded $H^1\to L^2$ and $M$ an isomorphism of $H^m$: so $E^{\rm add}$ is bounded iff $E'$ is. If $E'$ is unbounded in $H^{m+1}$, then (applying this with $m+1\ge1$) $E^{\rm add}$ has $\|E^{\rm add}u_j\|/\|u_j\|_{H^{m+1}}\to\infty$, and $\|u_j\|_{H^m}\le\|u_j\|_{H^{m+1}}$ gives the same blow-up in $H^m$. For $N=1$, P1-T2(i) gives $e^{sa_1(\partial)}v_0=v(t)$, so $E'=0$ and $E^{\rm add}\circ M=B$; for $N=0$, $E'(v)=v(t)-v_0$ is bounded $H^1\to L^2$ as in (a).

\step{3}{(c) Monomial terms}
If at least two monomials have degree $\ge1$, the no-go of Chapter~\ref{ch:cI3} (OPA-T3(b), unbounded for every $m$ and $t_0$) applies to $Q'$, and (b) transfers it to $Q$. With at most one, (b) gives boundedness. For reaction--diffusion, (a) with $N=1$ gives $E'=0$ and the formula; the mode computation is the product $(e^{sD\zeta^2}-1)(e^{sr}-1)$ of the two single-term increments (A19).

\whyworks{boundedness of the interaction is decided by the non-constant part of the operator; constants only multiply by a bounded exponential.}
\end{proof}
\uses{A19, A21, A22; P1-T2(i), OPA-T3 (Chapter~\ref{ch:cI3}).}

\begin{example}{The numerical dichotomy}
On $\Omega=(0,1)^2$ in $H^1$: for reaction--diffusion the ratio $R(\zeta)$ decreases, $0.031,0.0084,0.0021,0.00054$ for $\zeta=5,10,20,40$; for advection--diffusion--reaction $u_t=u_{xx}-u_x+2u$ it explodes, $0.79,\,30.9,\,1.7\cdot10^5,\,2.0\cdot10^{13}$, with the same dichotomy for $t_0=\tfrac12$ and $t_0=1$.
\end{example}

\section{Holes: the sum of kernels is closed but not dense}

\begin{dimcard}{A53}{S10A-P1}{On a region with a hole the sum of kernels is not dense (P2-Q1, C12-P7)}{\proved}{prop:s10ap1}
Let $P=\Delta=4\partial_z\partial_{\bar z}$, $Q_1=\partial_t+i\partial_x$, $Q_2=\partial_t-i\partial_x$ (coprime, $\delta(P)=1$); $\mathcal H_1$ are the holomorphic and $\mathcal H_2$ the anti-holomorphic functions in $H^m$.
\begin{enumerate}[label=(\roman*)]
\item On the annulus $A=\{1<|z|<2\}$ (Lipschitz, not convex, $b_1=1$), for every $m\ge0$, $\log|z|\in\ker_{H^m(A)}\Delta\setminus\overline{\mathcal H_1+\mathcal H_2}$: the sum is \emph{not dense}. For every $m\ge1$, moreover,
\[
\ker_{H^m(A)}\Delta=(\mathcal H_1+\mathcal H_2)\oplus\C\log|z| :
\]
the sum is \emph{closed, of codimension $1=b_1$, and not dense}.
\item If $\Omega$ is bounded Lipschitz with a cycle $\gamma\subset\Omega$ and a point $a\notin\Omega$ with $\mathrm{Ind}_\gamma(a)\ne0$, then $\log|z-a|\notin\overline{\mathcal H_1+\mathcal H_2}$ for every $m\ge0$.
\end{enumerate}
Consequently P2-Q1 is \textbf{false} on non-convex regions in general (density holds on convex regions, A18), and (i) is a \emph{differential} example of ``closed but not dense'', the question left open in C12-P7.
\end{dimcard}

\begin{proofidea}
A flux through a closed curve is a continuous functional on harmonic functions in $L^2$; it vanishes on holomorphic and anti-holomorphic functions (their derivative has a single-valued primitive) but not on $\log|z|$.
\end{proofidea}

\begin{proof}
\step{1}{The flux functional is continuous}
For a smooth cycle $\gamma\Subset\Omega$ and $h\in\ker_{L^2}\Delta$ put $\Phi_\gamma(h)=\int_\gamma\partial_nh\dd s=\int_\gamma{*dh}$. With $\delta=\mathrm{dist}(\gamma,\partial\Omega)/2$ the mean-value property and Cauchy's gradient estimate on a ball of radius $\delta/2$ give, for $y\in\gamma$,
\[
|\nabla h(y)|\le C'\delta^{-2}\|h\|_{L^2(\Omega)},\qquad\text{hence}\qquad |\Phi_\gamma(h)|\le C_\gamma\|h\|_{L^2}\le C_\gamma\|h\|_{H^m}.
\]
(Numerically on the annulus, $\sup_{r=3/2}|\partial_rh|/\|h\|_{L^2}\le0.78$ on $r^{\pm n}\cos n\theta$, $n\le8$.)

\step{2}{It vanishes on $\mathcal H_1+\mathcal H_2$ and not on $\log|z|$}
For $f$ holomorphic, $df=f'dz$ and ${*dz}=-i\,dz$, so ${*df}=-if'dz$ and $\Phi_\gamma(f)=-i\oint_\gamma f'(z)\dd z=0$ because $f'$ has the single-valued primitive $f$; the anti-holomorphic case is conjugate. For $h=\log|z|$ on a circle, $\Phi(h)=2\pi\ne0$. By continuity $\Phi_\gamma$ vanishes on the closure of $\mathcal H_1+\mathcal H_2$, so $\log|z|$ is not in it. This proves (i)'s first claim for every $m\ge0$, and (ii) the same way with the cycle $\gamma$ and the point $a$.

\step{3}{Exact codimension on the annulus for $m\ge1$}
A harmonic $h$ on $A$ has a Laurent form $h=f+\bar g+c\log|z|$ with $f,g$ holomorphic \emph{single-valued} and $c=\Phi(h)/2\pi$. If $\Phi(h)=0$ then $h=f+\bar g$, and one must check $f,\bar g\in H^m(A)$. Now $f'=\partial_zh\in H^{m-1}(A)$ (uses $m\ge1$). In polar coordinates, for $1<r<2$ and $r_0=3/2$,
\[
|f(re^{i\theta})-f(r_0e^{i\theta})|^2\le|r-r_0|\int_1^2|f'(\varrho e^{i\theta})|^2\dd\varrho ,
\]
so $\|f-f(r_0e^{i\cdot})\|_{L^2(A)}\le C\|f'\|_{L^2(A)}$ and $f(r_0e^{i\theta})$ is bounded; the derivatives of $f$ of order $1..m$ are derivatives of $f'\in H^{m-1}$, hence $f\in H^m(A)$; likewise for $\bar g$. So $\ker\Phi=\mathcal H_1+\mathcal H_2$ and $\ker_{H^m}\Delta=\ker\Phi\oplus\C\log|z|$, with $\ker\Phi$ closed (kernel of a continuous functional).

\whyworks{the obstruction is topological: a hole gives a period, and a period is a continuous functional that the sum of holomorphic and anti-holomorphic pieces cannot see.}
\end{proof}
\uses{A13, A14, A17, A18.}

\begin{caution}[title=What remains open]
(1) Non-Lipschitz regions (for example a disc with a slit, simply connected but not Lipschitz). (2) The exact codimension for general $P$ on a multiply connected region: for the Laplacian the flux gives it; in general one would need continuous ``period'' functionals in $H^m$. (3) Closedness on the annulus for $m=0$. In Phase~10 the Lipschitz simply connected case was proved modulo two cited facts [F1], [F2] (S10A-T3, A18 extension), and the extension of A18 to non-convex regions was shown to be \emph{false} (S10B-C1).
\end{caution}
